% Paper 1: Exit Dominance in Monetary Coordination Games
% Clean rebuild from SSRN preprint (6299081, Feb 2026) + Ledger formatting / citation pass
% + Theorem 1 ratchet form
% Build: xelatex BGT-PAPER-1.tex

\documentclass[12pt]{article}

% --- Fonts and encoding ---
\usepackage{fontspec}
\setmainfont{Latin Modern Roman}
\usepackage{amsmath}
% Note: amssymb omitted---unicode-math provides all symbols for XeLaTeX
\usepackage{unicode-math}
\setmathfont{Latin Modern Math}

% --- Unicode glyph fallbacks (mandatory for TL2023+) ---
\usepackage{newunicodechar}
\newunicodechar{δ}{$\delta$}
\newunicodechar{ε}{$\varepsilon$}
\newunicodechar{λ}{$\lambda$}
\newunicodechar{σ}{$\sigma$}
\newunicodechar{κ}{$\kappa$}
\newunicodechar{ρ}{$\rho$}
\newunicodechar{Δ}{$\Delta$}
\newunicodechar{Σ}{$\Sigma$}
\newunicodechar{∞}{$\infty$}

% --- Page layout ---
\usepackage[margin=1in]{geometry}
\usepackage{setspace}
\onehalfspacing

% --- Theorem environments ---
\usepackage{amsthm}
\newtheorem{theorem}{Theorem}
\newtheorem{proposition}{Proposition}
\newtheorem{lemma}{Lemma}
\newtheorem{corollary}{Corollary}
\theoremstyle{definition}
\newtheorem{definition}{Definition}
\newtheorem{assumption}{Assumption}
\theoremstyle{remark}
\newtheorem*{remark}{Remark}

% --- Tables ---
\usepackage{booktabs}
\usepackage{longtable}
\usepackage{array}
\usepackage{tabularx}
\usepackage{caption}

% --- Page numbers ---
\usepackage{fancyhdr}
\pagestyle{fancy}
\fancyhf{}
\fancyfoot[C]{\thepage}
\renewcommand{\headrulewidth}{0pt}

% --- Misc ---
\usepackage{enumitem}
\usepackage{hyperref}
\hypersetup{
  colorlinks=true, linkcolor=black, citecolor=black, urlcolor=blue,
  pdftitle={Bitcoin Exit Dominance in Monetary Coordination Games},
  pdfauthor={Sean Hash},
  pdfkeywords={Bitcoin, game theory, Nash equilibrium, monetary coordination, neutral settlement, capital-weighted adoption, path dependence, coordination failure, threshold strategy, monotone comparative statics, digital scarcity, Schelling point, network effects}
}
\def\UrlBreaks{\do\/\do-\do.\do=\do_\do?\do\&\do\%\do\a\do\b\do\c\do\d\do\e\do\f\do\g\do\h\do\i\do\j\do\k\do\l\do\m\do\n\do\o\do\p\do\q\do\r\do\s\do\t\do\u\do\v\do\w\do\x\do\y\do\z\do\A\do\B\do\C\do\D\do\E\do\F\do\G\do\H\do\I\do\J\do\K\do\L\do\M\do\N\do\O\do\P\do\Q\do\R\do\S\do\T\do\U\do\V\do\W\do\X\do\Y\do\Z}
\usepackage{microtype}

% --- Suppress section numbering (manual numbers in headings) ---
\setcounter{secnumdepth}{0}

% --- Custom commands ---
\newcommand{\E}{\mathbb{E}}

\begin{document}

% ============================================================
% TITLE BLOCK
% ============================================================

\begin{center}
{\LARGE\bfseries Bitcoin Exit Dominance in Monetary Coordination Games}

\bigskip

{\large Sean Hash}\\
\url{bitcoingametheory.com}\\
\href{mailto:sean@bitcoingametheory.com}{sean@bitcoingametheory.com}

\end{center}

\bigskip

% ============================================================
% ABSTRACT
% ============================================================

\begin{abstract}
\noindent
\textbf{Abstract.} In a multipolar world with no trusted monetary coordinator,
how do rational actors settle large-value transactions across trust
boundaries? This paper models the choice as a non-cooperative game (the ``Exit Game'') in
which capital allocators choose between capturable settlement systems
(``Stay'') and neutral settlement (``Exit''). The model rests on four assumptions: persistent multipolarity, rational
self-interest, computational hardness, and the persistence of network effects. Three results follow. First, the payoff advantage of Exit over Stay
is strictly increasing in adoption: every payoff term favors Exit,
and actors holding meaningful capital see their tipping points fall
under structural debasement (Proposition~1). Second, no Stay coalition
can hold itself together once any member's tipping point has been
crossed under open access (Proposition~2). Third, above the critical mass $p_c$ no Stay coalition can re-form, and the set of starting points from which adoption reaches the high-Exit configuration grows over time (Theorem~1). Together these establish Bitcoin Exit Dominance (Theorem~2).
\end{abstract}

\newpage

% ============================================================
% 1. INTRODUCTION
% ============================================================

\section{1. Introduction}

A capital allocator with purchasing power to protect faces a choice between two settlement systems. The first is the prevailing one: sovereign currency, custodied by regulated intermediaries, settled through the banking and securities infrastructure. Call this Stay. The second is a settlement asset that no government issues, no intermediary custodies, and no political body can debase or freeze. Call this Exit. The interesting question is not which a single allocator prefers, but what happens when every allocator faces the same choice at the same time.

The choice became possible on 3 January 2009. Nakamoto\textsuperscript{21}'s Bitcoin client mined its first block, and embedded in the coinbase transaction was the front-page headline of \emph{The Times} of London that morning: ``Chancellor on brink of second bailout for banks.'' The block carried no financial information. It carried a pointer to what the proposed system was meant to escape: the bailout regime that absorbs private failures into public balance sheets and pays for them through monetary debasement. The genesis block is the question this paper formalizes. When capturable settlement systems can be debased without recourse, can rational actors settle outside them?

The question is structural, not conjunctural. Global public debt approaches 100\% of GDP, with adverse scenarios projecting roughly 115\% by the early 2030s (IMF\textsuperscript{15}); historical repayment of sovereign debt at this scale has occurred through financial repression and debasement rather than taxation (Reinhart and Sbrancia\textsuperscript{23}). What allocators do against that backdrop, individually and in aggregate, determines whether neutral settlement scales or stays a curiosity.

This paper takes the choice and makes it a game. Each allocator picks Exit or Stay observing the current capital-weighted Exit fraction $p_t$. Payoffs depend on $p_t$: as more capital exits, the Exit side becomes more attractive through deeper markets, lower volatility, and normalized regulation; the Stay side becomes worse through rising non-adoption penalties. The question is whether this game has an equilibrium in which $p_t$ stays low forever, oscillates between regimes, or admits a critical mass past which no group of actors has reason to switch back to Stay.

Seventeen years of empirical record point at the third equilibrium. Spot Bitcoin exchange-traded products cleared U.S.\ regulatory approval in January 2024 (SEC\textsuperscript{27}); BlackRock's iShares Bitcoin Trust accumulated roughly \$37 billion in cost basis within its first fiscal year (BlackRock\textsuperscript{5}); central bank gold reserves rose at the fastest pace since the 1970s after the 2022 sanctions on Russian foreign-exchange holdings (Arslanalp, Eichengreen, and Simpson-Bell\textsuperscript{1}). This paper formalizes the game already in motion.

\newpage

% ============================================================
% 2. RELATED WORK
% ============================================================

\section{2. Related Work}

The closest prior work is Brunnermeier, James, and
Landau\textsuperscript{6}, who pose the same coordination problem.
Their response --- institutional coordination via central bank
digital currencies --- differs from this paper's premise
(Assumption~1) and fails its test: a CBDC digitizes the capturable
system rather than escape it, satisfies neither D1 nor a load-bearing
subset of P1--P7, and is an instance of Stay on which Theorem~2
operates.

The existing Bitcoin game-theory literature analyzes a different
layer of the same system. Biais, Bisiere, Bouvard, and
Casamatta\textsuperscript{4} show that following the longest chain
is an equilibrium of the miner game; Eyal and
Sirer\textsuperscript{9} characterize selfish mining above a
one-third hashrate threshold; Budish\textsuperscript{7} bounds
attack cost in network value; and Chen\textsuperscript{8} extends
the security-budget argument to fee markets. All four study the
miner game. The present paper analyzes the allocator meta-game one
layer above: capital holders choose between capturable and neutral
settlement. The miner game is proved in prior work and not re-derived
here.

% ============================================================
% 3. MODEL
% ============================================================

\section{3. Model}

\subsection{3.1 Maintained Assumptions}

The analysis rests on four assumptions. Each is independently assessable. All subsequent claims derive from these assumptions; rejecting any assumption invalidates the specific claims that depend on it.

\begin{assumption}[Multipolarity]
No single entity permanently governs global economic activity. Power
distributes across competing centers with superlinear coordination
costs. The multipolar configuration is dynamic. The institutional-fragmentation index $D_t$ (\S3.6) is monotone non-decreasing on the long-run trajectories relevant to Theorem~1.
\end{assumption}

\begin{assumption}[Rational Self-Interest]
Actors optimize for self-interest. When defection from cooperative
agreements is unpunished or unpunishable, actors defect.
\end{assumption}

\begin{assumption}[Computational Hardness]
Certain mathematical problems remain computationally intractable.
Digital scarcity and cryptographic custody are implementable.
\end{assumption}

\begin{assumption}[Network Effect Persistence]
Past critical mass, switching costs exceed marginal gains of
alternatives. Incumbency compounds (Katz and
Shapiro\textsuperscript{16}).
\end{assumption}

\subsection{3.2 Definitions}

\begin{definition}[Neutral Settlement]
An asset $S$ is a \emph{neutral settlement asset} if and only if $S$ is immune to capture by force, debasement, or coordinated political action. Table~1 lists the seven properties P1--P7 jointly necessary for D1, alongside the attack class each blocks. Property P6 (neutrality) corresponds to what Huberman, Leshno, and Moallemi\textsuperscript{14} term ``monopoly without a monopolist'': the protocol exhibits monopoly-like network effects without an entity capable of exercising monopoly power. The cryptographic and consensus primitives that ground P1--P3 and P5--P6 trace to Nakamoto\textsuperscript{21}. Bitcoin satisfies P1--P7. The formal definition, irreducibility, and uniqueness of P1--P7 across asset classes are taken up in a companion paper in preparation; the present paper takes P1--P7 as maintained assumptions (\S3.2), as it does the miner sub-game result.
\end{definition}

\begin{table}[h]
\centering
\caption*{\textbf{Table 1.} The seven properties jointly necessary for Definition~1. Each row names what an asset must guarantee and the attack class it forecloses.}
\begin{tabular}{ll}
\toprule
Property & Blocks \\
\midrule
P1: Informational security & Physical confiscation \\
P2: Protocol security & Protocol rule change \\
P3: Permissionless access & Transaction censorship \\
P4: Absolute scarcity & Supply inflation \\
P5: Cheap finality & Prohibitive cost/delay \\
P6: Neutrality & Governance takeover \\
P7: Adaptive resilience & Obsolescence \\
\bottomrule
\end{tabular}
\end{table}

Gold is the historical neutral-settlement instrument and remains relevant: central bank gold reserves have risen since the Global Financial Crisis, accelerating in correlation with financial-sanctions risk on the main reserve issuers (Arslanalp, Eichengreen, and Simpson-Bell\textsuperscript{1}). Density, divisibility, recognizability, and aesthetic appeal have sustained gold as money for five millennia in the market-selection account of monetary emergence (Menger\textsuperscript{19}). Other proto-monies including wampum, cowrie shells, and salt shared these surface properties but failed P4: industrial production, oceanic shipping, and refined extraction each broke their scarcity. Only gold has held P4 across millennia of technological advance. Gold's known structural failures are physical seizure, with the 1933 Executive Order 6102 as the canonical precedent; assay risk at unit level, exemplified by tungsten-cored bar fraud; and settlement cost that scales with physical mass. Bitcoin closes these failures while preserving gold's monetary properties; the formal P1--P7 treatment is taken up in the companion paper in preparation.


\begin{definition}[Capturable System]
A monetary system in which at least one of P1--P7 fails for its
underlying settlement asset.
\end{definition}

\begin{definition}[Exit]
The action of moving capital from a capturable system to a neutral
settlement asset; distinct from Hirschman's\textsuperscript{13}
exit/voice/loyalty triad, which describes responses to
organizational decline. The complement is \emph{Stay}: holding
capital in the prevailing capturable system.
\end{definition}

\begin{definition}[Critical Mass]
The critical-mass threshold at slowly varying state $(K, D)$ is
\[
p_c(K, D) := \inf\left\{ p \in [0, 1] : \sum_{i\,:\,\Delta_i(p; K, D) > 0} a_i \;>\; 0 \right\},
\]
where $\Delta_i$, $K$, and $D$ are defined in \S3.6.
\end{definition}


\subsection{3.3 Game}

Monetary coordination is modeled here as a non-cooperative game. Each capital allocator chooses one of two
actions: hold capital in the prevailing capturable system (Stay),
or move it to a neutral settlement asset outside that system
(Exit).

\smallskip
\noindent\textbf{Players.} The set $N$ is the population of
capital allocators, weighted by their share $a_i$ of total
capturable capital. Sovereigns, central banks, corporates, asset
managers, and household wealth all participate. Headcount does
not move the game; capital does.

\smallskip
\noindent\textbf{Objective.} Each player maximizes utility. The
Exit side nets real return against volatility, switching cost,
and regulatory exposure. The Stay side nets capturable return
against capturable volatility and a non-adoption penalty. Players
differ in how much they weight each cost, captured by
heterogeneity parameters $\lambda_i, \kappa_i, \rho_i$.

\smallskip
\noindent\textbf{Rules.} Each player observes the current
capital-weighted adoption fraction $p_t \in [0, 1]$ and picks
Exit or Stay. As more capital exits, the
Exit side becomes more attractive: deeper markets, lower
volatility, regulation normalizes. The Stay side weakens through
rising non-adoption penalties. Exact payoffs are stated
in \S3.4; the monotonicity conditions governing how they respond
to $p$ are stated in \S3.5.

\smallskip
\noindent\textbf{Dynamics.} Background conditions evolve slowly.
Recognition of the financial-repression channel spreads across capital cohorts, and
institutional fragmentation trends upward across multi-decade horizons.
Both push individual thresholds down, gradually moving the game
toward the critical mass past which Stay coalitions can no longer hold together
(\S3.6). Section~4 says what makes ignition possible, what blocks ignition,
and why Exit, once won, does not reverse.

\subsection{3.4 Payoffs}

The game is $G = (N, S, u)$ in the standard non-cooperative form (Fudenberg and Tirole\textsuperscript{12}). The player set $N = \{1, \ldots, n\}$ is the population of capital allocators, where actor $i$ holds capital share $a_i \in [0,1]$ with $\sum_i a_i = 1$. Each strategy set is $S_i = \{\text{Exit}, \text{Stay}\}$. The aggregate adoption process is $p_t = \sum_{i \in \text{Exit}_t} a_i \in [0,1]$, the capital-weighted fraction at Exit. This dollar-weighted convention follows monetary economics, where reserve composition and capital flows are weighted by stake rather than by headcount: holdouts with negligible capital weight do not move $p_t$.

Because payoffs depend on others only through the aggregate $p$, the utility function $u_i : S \to \mathbb{R}$ is written in two cases:
\[
u_i^E(p) := R_B(p) - \lambda_i \sigma_B(p) - \kappa_i K_A(p) - \rho_i R_A(p),
\]
\[
u_i^S(p) := R_F - \lambda_i \sigma_F - K_N(p).
\]

\smallskip
\noindent\textbf{Exit payoff components.}
\begin{itemize}[nosep,leftmargin=*]
\item $R_B(p)$: real return on the neutral asset.
\item $\lambda_i \sigma_B(p)$: volatility cost, scaled by personal risk aversion.
\item $\kappa_i K_A(p)$: adoption penalty, the switching cost of moving to neutral settlement (Katz and Shapiro\textsuperscript{16}; Farrell and Saloner\textsuperscript{10}), scaled by personal career-risk loading.
\item $\rho_i R_A(p)$: regulatory penalty, scaled by personal regulatory exposure.
\end{itemize}

\smallskip
\noindent\textbf{Stay payoff components.}
\begin{itemize}[nosep,leftmargin=*]
\item $R_F$: capturable return.
\item $\lambda_i \sigma_F$: volatility cost on the capturable side.
\item $K_N(p)$: non-adoption penalty, a negative network externality on holdouts as competitors exit.
\end{itemize}

\smallskip
\noindent Each actor carries three personal weights: $\lambda_i$ on
volatility, $\kappa_i$ on the career risk of switching infrastructure,
and $\rho_i$ on regulatory exposure. All three are strictly positive
and differ across actors, which is why the threshold $p_i^*$ at
which Exit becomes attractive varies across the population: a
sovereign wealth fund and a retail allocator face the same payoff
equation but cross at different $p$. The functions $R_B$, $\sigma_B$, $K_A$, $R_A$, $K_N$ are continuous and bounded on $[0,1]$. $\sigma_B(p)$ is a quadratic risk measure consistent with mean-variance preferences. Bitcoin's expected return $R_B(p)$ and volatility
$\sigma_B(p)$ both depend on adoption because adoption raises
the dollar value of the block reward that pays for protocol
security (Pagnotta\textsuperscript{22}; Chen\textsuperscript{8}).

\smallskip
\noindent Table~2 collects $u_i^E$ and $u_i^S$ in normal form: the
row is player $i$'s action, the column is what the rest of the
capital is doing, and each cell is the payoff from the equations
above evaluated at the corresponding $p$.

\begin{table}[h]
\centering
\caption*{\textbf{Table 2.} Normal-form payoff matrix for player $i$, in the shorthand $u_i^E$, $u_i^S$ defined above.}
\begin{tabular}{lcc}
\toprule
 & Others: Stay & Others: Exit \\
\midrule
Player $i$: Stay & $u_i^S(0)$ & $u_i^S(p)$ \\
Player $i$: Exit & $u_i^E(0)$ & $u_i^E(p)$ \\
\bottomrule
\end{tabular}
\end{table}

\subsection{3.5 Monotonicity Conditions}

Five conditions govern how the payoff primitives respond to rising adoption $p$.

\begin{enumerate}[label=(M\arabic*),nosep]
\item $R_B'(p) > 0$: Network effects increase return (Assumption~4)
\item $\sigma_B'(p) < 0$: Deeper markets reduce volatility
\item $K_A'(p) < 0$: Adoption penalty falls with adoption
\item $R_A'(p) < 0$: Regulatory penalty decreases as adoption normalizes compliance
\item $K_N'(p) > 0$: Non-adoption penalty rises as competitors exit
\end{enumerate}

\smallskip
\noindent Each condition specifies a single-direction response to
rising adoption. Three concern the
Exit side becoming more attractive as adoption deepens: returns
rise (M1), volatility falls (M2), and switching costs fall (M3).
One concerns regulatory friction normalizing as adoption becomes
mainstream (M4). One concerns the externality on non-adopters
strengthening as competitors exit (M5). These are empirical conditions, maintained to hold on average over multi-decade horizons. \S5 returns to this point.

\subsection{3.6 Threshold Drift}

The capturable system's failure mode is not the overt seizure,
capital control, or hyperinflation that history records as
discrete events. It is the continuous extraction of real wealth
through the interaction of nominal asset appreciation and the
taxation of nominal gains. As debasement raises the price level,
equities, real estate, and similar capturable assets appreciate
nominally to preserve real value. Holders are then taxed on the
nominal gain. Post-tax real returns are negative (Feldstein\textsuperscript{11}); the public-debt
numerator grows while the real value of private claims is
transferred to service it. The decline is imperceptible because
nominal balances do not contract. The loss registers only when
measured against the price level. This is the channel through which
Assumption~1 delivers $R_F < 0$ in real terms, and it is what
cohorts come to recognize over multi-decade horizons.

Let $K_t \in [0,1]$ be the share of capital cohorts that recognize
capturable assets as financially repressed rather than risk-free. Recognition need not arrive through a discrete capture event but accumulates as cohorts observe that nominal returns, after tax on nominal capital gains, fail to keep pace with the price level. Let $D_t \in [0, \infty)$ be an
institutional-fragmentation index summarizing four inputs: debt-to-GDP,
the inflation--capital-gains tax wedge, financial-repression incidence,
and currency-block fragmentation. Both are slowly varying
state variables, monotone non-decreasing on the multi-decade
horizons relevant to Theorem~1, as observed in post-1970 fiat
regime trends (Reinhart and Sbrancia\textsuperscript{23}).

Each actor's threshold is a function of the slowly varying state:
\[
p_i^*(K, D) := \inf\{p \in [0, 1] : \Delta_i(p; K, D) > 0\}.
\]
Assume the sign conditions
\[
\frac{\partial p_i^*}{\partial K} \leq 0, \qquad
\frac{\partial p_i^*}{\partial D} \leq 0.
\]
Strict inequality holds on a set of positive capital weight. Higher $D_t$ lowers $R_F$ in real terms. Higher $K_t$ erodes $\kappa_i K_A$ and $\rho_i R_A$. Fewer actors still treat the capturable system as a safe numeraire. Both
effects raise $\Delta_i(0; K, D)$ and push $p_i^*$ down.

The critical mass $p_c(K, D)$ defined in \S3.2 inherits the same
sign structure: rising $D_t$ and rising $K_t$ push $p_c$ down over multi-decade horizons, lowering
the bar for ignition. Below $p_c$ no positive-capital-weight set
prefers Exit; at or above $p_c$, defection by some positive-weight
set is incentivized. Whether adoption continues or stalls depends on the threshold density of the actor population. If successive thresholds are sparse relative to first-defector capital weights, the cascade may pause. Further drift in $(K_t, D_t)$ resumes ignition. This is
the path-dependent adoption mechanism of Arthur\textsuperscript{2}:
small differences in early allocation history, compounded by
increasing returns to adoption, select the long-run regime.

The model's claims are structural and directional: existence of
thresholds $\{p_i^*\}$, sign of comparative statics, and
direction of basin expansion under accumulating $(K_t, D_t)$. It
does not predict specific threshold values, tipping times, or
saturation levels, which depend on functional forms and
parameter calibrations the model leaves unspecified. This
follows the qualitative comparative-statics paradigm of Milgrom
and Shannon\textsuperscript{20} and Topkis\textsuperscript{26}
and the path-dependent lock-in paradigm of
Arthur\textsuperscript{2}.

An adoption level $p$ is an equilibrium of the Exit Game when no actor, alone or in any group that holds together, gains by switching. \S4 establishes that above $p_c$ this holds: the high-Exit outcome survives any returning coalition, and the set of starting points that reach it only grows as $(K_t, D_t)$ accumulate.


% ============================================================
% 4. RESULTS
% ============================================================

\section{4. Results}

The argument proceeds in four steps. \S4.1 fixes the static
incentive: the per-period advantage of Exit over Stay rises with
adoption (Proposition~1). \S4.2 takes that incentive into the
coordination problem: no coalition holding Stay survives unilateral
defection (Proposition~2). \S4.3 takes the coordination result into
the dynamic problem: above the critical mass $p_c$ the high-Exit configuration holds against any defecting group, and under threshold drift
more starting points lead to the high-Exit configuration (Theorem~1). \S4.4 composes the three into the synthesis
claim, Bitcoin Exit Dominance (Theorem~2). Each step uses the previous;
none of the three intermediate results is the title claim on its
own.

\subsection{4.1 Payoff Advantage}

Hold the slowly varying state $(K, D)$ fixed. The payoff differential $\Delta_i(p)$ rises strictly with $p$ for every player type. Each player has a threshold $p_i^*$ above which Exit is the unique best response. By convention, $p_i^* = +\infty$ for permanent holdouts whose $\Delta_i$ stays non-positive on $[0, 1]$.

\begin{proposition}[Monotone Exit Advantage]
Under Assumptions~1--4 and maintained conditions (M1)--(M5):

\begin{enumerate}[label=(\roman*)]
\item The payoff differential $\Delta_i(p)$ is strictly increasing
in $p$ for all player types.
\item For each player $i$ there exists a threshold
$p_i^* \in [0, 1] \cup \{+\infty\}$ such that Exit is the unique
best response for $p > p_i^*$.
\end{enumerate}
\end{proposition}

\begin{proof}
Define the payoff differential:
\[
\Delta_i(p) = u_i(\text{Exit}) - u_i(\text{Stay})
\]
\[
= [R_B(p) - R_F] - \lambda_i[\sigma_B(p) - \sigma_F] - \kappa_i \cdot K_A(p) - \rho_i \cdot R_A(p) + K_N(p).
\]

Taking the derivative with respect to $p$:
\[
\frac{d\Delta_i}{dp} = R_B'(p) - \lambda_i \cdot \sigma_B'(p) - \kappa_i \cdot K_A'(p) - \rho_i \cdot R_A'(p) + K_N'(p).
\]

Under (M1)--(M5), each term is positive (Table~3).

\begin{table}[h]
\centering
\caption*{\textbf{Table 3.} Sign of each term of $d\Delta_i/dp$ under (M1)--(M5).}
\begin{tabular}{lcl}
\toprule
Term & Sign & Reason \\
\midrule
$R_B'(p)$ & $> 0$ & (M1) \\
$-\lambda_i \cdot \sigma_B'(p)$ & $> 0$ & $\lambda_i > 0$, $\sigma_B'(p) < 0$ by (M2) \\
$-\kappa_i \cdot K_A'(p)$ & $> 0$ & $\kappa_i \geq 0$, $K_A'(p) < 0$ by (M3) \\
$-\rho_i \cdot R_A'(p)$ & $> 0$ & $\rho_i \geq 0$, $R_A'(p) < 0$ by (M4) \\
$K_N'(p)$ & $> 0$ & (M5) \\
\bottomrule
\end{tabular}
\end{table}

\noindent Therefore $d\Delta_i/dp > 0$: the advantage of Exit rises strictly
with $p$. By the monotone comparative-statics results of Milgrom
and Shannon\textsuperscript{20} and Topkis\textsuperscript{26}, each
actor's best response is non-decreasing in $p$.

Define each player's threshold as
\[
p_i^* := \inf\{p \in [0, 1] : \Delta_i(p; K, D) > 0\},
\]
with the convention $p_i^* = +\infty$ when the set is empty. For
actors with $p_i^* \in [0, 1]$, the intermediate value theorem
applied to the continuous $\Delta_i$ delivers
$\Delta_i(p_i^*) = 0$, and for $p > p_i^*$ Exit strictly dominates.
Actors with $p_i^* = +\infty$ never defect within the unit
interval; the high-Exit region Theorem~1 establishes
admits them by construction.

Under Assumption~1, $R_F < 0$ in real terms (see \S1; Reinhart and Sbrancia\textsuperscript{23}). This ensures that even at $p = 0$, the
term $[R_B(0) - R_F]$ provides positive contribution, lowering
$p_i^*$ toward zero for the actors whose remaining parameters are
not extreme.

\end{proof}

\smallskip
\noindent When $R_F < 0$, the \S3.4 payoff matrix (Table~2) collapses: for any actor whose threshold $p_i^*$ has been crossed, Exit pays more than Stay regardless of what others do.

\subsection{4.2 Coordination Failure}

Proposition~1 gives the per-actor incentive but says nothing about
group behavior. A coalition might still try to hold Stay
collectively. It cannot. Standard free-rider logic disqualifies
any all-Stay coalition with crossed-threshold members from being
self-enforcing.

\begin{proposition}[Coalition Instability]
Under Assumption~2 and Property~P3, no Stay coalition $C$
containing at least one actor with a crossed threshold
($p > p_j^*$ for some $j \in C$) is self-enforcing. The result
is in the spirit of Bernheim-Peleg-Whinston self-enforcement,
applied here at the unilateral level: when defection is
unpunishable, the crossed-threshold member defects, breaking the
coalition at the Nash layer that any CPNE refinement requires.
\end{proposition}

\begin{proof}
Consider a coalition $C \subseteq N$ maintaining Stay. For this
coalition to be self-enforcing in the sense of Bernheim, Peleg, and
Whinston\textsuperscript{3}, no member $j \in C$ can profit from
unilateral defection.

But by Proposition~1, $u_j(\text{Exit}) > u_j(\text{Stay})$ for $j$
when $p$ exceeds $p_j^*$. Property P3 (open access)
makes the institutional defection cost zero: no counterparty can
refuse the trade and no clearing intermediary can foreclose
settlement. Sanctions enter the payoff function via $R_A$ and are already priced into each actor's threshold; Proposition~2's claim presumes that pricing.

The coalition faces a standard free-rider problem. The benefit of
collective Stay is a public good (maintained value of capturable
system); the benefit of individual defection is a private good
(captured network effects, avoidance of debasement). Under
Assumption~2, the private good dominates.

The dynamics compound: the first mover captures fleeing
capital, raising $K_N(p)$ for remaining coalition members and
accelerating further defection.

Formally, for at least one $j \in C$,
\[
u_j(\text{defect from } C) > u_j(\text{remain in } C),
\]
which is sufficient to violate Bernheim-Peleg-Whinston
coalition-proofness.

The first-mover advantage is empirically observable: El Salvador adopted Bitcoin as legal tender in 2021 (Republic of El Salvador\textsuperscript{24}), and the U.S. Securities and Exchange Commission approved spot Bitcoin exchange-traded products in January 2024 (SEC\textsuperscript{27}), with BlackRock's iShares Bitcoin Trust accumulating roughly \$37 billion in cost basis within its first fiscal year (BlackRock\textsuperscript{5}).
\end{proof}

\begin{remark}
Coalitions of permanent-holdout actors ($p_j^* = +\infty$ for all
$j \in C$) are admitted by Proposition~1's threshold convention
and consistent with the high-Exit region Theorem~1 establishes.
Such coalitions, when maintained by coercive sanction against P3,
are subject to preference falsification
(Kuran\textsuperscript{18}): publicly observed Stay coexists with
privately optimal Exit, and the first-mover advantage operates on
private allocations regardless of public posture. P3 reduces the technical cost of private holding to zero. A dimensionless social cost remains for maintaining the public Stay posture while privately exiting. Under the multipolar enforcement of Assumption~1 (no global authority, jurisdictional arbitrage, imperfect detection), this social cost is bounded for any individual actor. The capital-preservation benefit under $R_F < 0$ is not bounded. Once
accumulated private Exits show up in $p_t$, the public Stay
posture no longer holds for crossed-threshold actors.
\end{remark}

\subsection{4.3 Ratchet}

This section establishes forward invariance of the absorbing region $\{p > p_c\}$ and monotone basin expansion under \S3.6 drift. The argument proceeds by case-split on candidate Stay coalitions: those intersecting the crossed-threshold set $H(p)$ break by Proposition~2; those disjoint from $H(p)$ do not grow.

\begin{theorem}[Ratchet Form; cf. Arthur 1989 lock-in\textsuperscript{2}]
Under Assumptions~1--4, maintained conditions (M1)--(M5), and
the \S3.6 sign conditions: (a) the absorbing region $\{p > p_c\}$ is forward invariant under perturbations remaining within it. (b) Under
monotonic accumulation of $(K_t, D_t)$, the critical mass
$p_c(K_t, D_t)$ drifts down, so the basin of local stability
expands monotonically over time. The ratchet does not reverse: once $p$ enters the region above $p_c$, the basin only grows.
\end{theorem}

\begin{proof}
\emph{(a) Local stability above $p_c$.} If $p > p_c(K, D)$, the set $H(p) = \{i : \Delta_i(p; K, D) > 0\}$ has
positive capital weight. Any perturbation attempting to coordinate
a return to high-Stay above $p_c$ is a candidate Stay coalition
$C$ and falls into one of two cases. If $C \cap H(p)$ is
nonempty, Proposition~2 applies: the crossed-threshold member
defects under P3, and $C$ is not self-enforcing. If $C \cap H(p)$
is empty, $C$ consists entirely of uncrossed-threshold capital. These actors are already at their preferred play at the current adoption level and do not move it. Neither case grows.
The high-Exit configuration is locally stable. The static part is standard
Bernheim-Peleg-Whinston coalition-proofness; the dynamic part,
in which the finite-threshold Stay pool shrinks as
$(K_t, D_t)$ accumulate, is in the spirit of Konishi and
Ray\textsuperscript{17}. Drawdowns and cascade pauses along $p_t$ are admitted. A permanent stall would require two conditions to hold jointly: $p_t$ plateauing below $p_c$ indefinitely, and $p_c$ failing to fall. The second condition is ruled out by part (b).

\emph{(b) Basin drift.} The \S3.6 sign conditions give
$\partial p_i^*/\partial K \leq 0$ and $\partial p_i^*/\partial D
\leq 0$. The critical mass inherits these signs: $\partial p_c/
\partial K \leq 0$ and $\partial p_c/\partial D \leq 0$. Under
monotonic accumulation of $(K_t, D_t)$, $p_c(K_t, D_t)$ is
non-increasing in $t$. The basin of local stability $\{p : p >
p_c(K_t, D_t)\}$ therefore expands monotonically. The ratchet
structure is the lock-in by historical events of
Arthur\textsuperscript{2} applied to monetary adoption: not
pointwise irreversibility, but monotone expansion of the
locally stable absorbing region as $(K_t, D_t)$ accumulate.
\end{proof}

\smallskip
\noindent The high-Exit region above $p_c$ allows partial adoption, in the manner
of incumbent network goods such as USD reserve status, English
as lingua franca, or Google search dominance, which empirically
plateau at finite shares rather than universal adoption (Katz
and Shapiro\textsuperscript{16}; Farrell and
Saloner\textsuperscript{10}). Permanent holdouts are admitted by construction. Some Stay capital remains, parked at a
focal coordination level (Schelling\textsuperscript{25}) that no
returning group can unwind. Adoption settles short of universal Exit.

\subsection{4.4 Bitcoin Exit Dominance}

Proposition~1 gives the static incentive. Proposition~2 rules out Stay coalitions with crossed-threshold members. Theorem~1 gives local stability above $p_c$ with monotone basin expansion. The realized process exhibits all three at once.
That conjunction is the title claim.

\begin{theorem}[Bitcoin Exit Dominance]
Under the joint hypotheses of Propositions~1--2 and Theorem~1,
the realized capital-weighted adoption process satisfies all
three of:
\begin{enumerate}[label=(\roman*),nosep]
\item the payoff advantage of Exit over Stay is strictly increasing
in $p$ for every player type (Proposition~1);
\item no Stay coalition with at least one crossed-threshold
member is self-enforcing (Proposition~2);
\item the high-Exit configuration is locally stable above
$p_c$, and the basin of local stability expands monotonically
under threshold drift (Theorem~1).
\end{enumerate}
\end{theorem}

\begin{proof}
Proposition~1 gives (i) directly under (M1)--(M5). Proposition~2
gives (ii) under Assumption~2 and Property P3, with (i) supplying
the per-defector incentive that breaks any Stay coalition with
at least one crossed-threshold member. Theorem~1 gives (iii):
local stability of the high-Exit configuration above $p_c$ via
Proposition~2 applied to candidate Stay coalitions, and basin
expansion via the \S3.6 sign conditions on $p_c(K_t, D_t)$. The
conjunction of (i), (ii), and (iii) is the Bitcoin Exit Dominance claim:
strictly increasing payoff advantage in $p$, instability of
crossed-threshold Stay coalitions, and a locally stable high-
Exit absorbing region whose basin expands under threshold drift.
\end{proof}

\begin{remark}
By the implicit function theorem applied to
$\Delta_i(p_i^*) = 0$,
\[
\frac{dp_i^*}{d\lambda_i}
\;=\;
-\,\frac{\partial \Delta_i / \partial \lambda_i}{d \Delta_i / dp}
\;=\;
\frac{\sigma_B(p_i^*) - \sigma_F}{d \Delta_i / dp},
\]
which is positive when $\sigma_B(p_i^*) > \sigma_F$, the
empirically relevant case of a higher-volatility neutral asset
relative to the capturable side. Under that condition, more
risk-averse actors carry higher thresholds and cross them later;
because $d\Delta_i/dp > 0$, they cross once $p$ rises
sufficiently.
\end{remark}

% ============================================================
% 5. FALSIFICATION
% ============================================================

\section{5. Limitations}

The paper takes P1--P7 as maintained assumptions on Definition~1. Their formal characterization, irreducibility, and uniqueness across asset classes are the subject of a companion paper in preparation; the present paper does not establish that no alternative asset satisfies P1--P7 jointly.

The model's claims are qualitative and directional: thresholds $\{p_i^*\}$ exist and fall as $(K_t, D_t)$ accumulate, and the set of starting points reaching the high-Exit region grows over time. It does not predict specific threshold values, tipping times, or saturation levels.

The monotonicity conditions (M1)--(M5) are stated to hold on average over multi-decade horizons rather than at every instant. Transient violations, for instance $R_B'(p) < 0$ during a network-congestion crisis, are within scope.

Adoption-timing phenomena are outside scope: price declines, regulatory actions against specific intermediaries, developer controversies, energy debates, and individual-cohort deleveraging affect when, not whether, the high-Exit region is reached.

\section{6. Conclusion}

The standard frame for monetary coordination treats sovereign issuers as the only agents that matter. The Exit Game inverts the frame: capital allocators decide what counts as settlement, and the result is a game whose long-run Exit outcome depends on the existence of an asset that no government issues. Once such an asset exists and satisfies P1--P7, the question is no longer whether allocators can leave the capturable system but whether continued debasement and fragmentation push adoption past $p_c$.

The implication for the international monetary system is that reserve composition depends on settlement-system characteristics rather than only on political coordination. Central bank gold accumulation since the Global Financial Crisis is the directional precedent. The Exit Game predicts that any settlement asset which closes gold's three structural failures of assay, confiscation, and supply growth pulls capital the same direction at lower cost; the \S3.6 threshold-drift mechanism predicts that continued financial repression and institutional fragmentation keep pushing more actors into Exit.

Two questions are deliberately deferred. Whether neutral settlement is unique to Bitcoin, or some other asset class can also satisfy P1--P7 jointly, is the subject of the companion paper in preparation. Whether the allocator population can be extended to include autonomous economic agents whose strategic structure differs from that of human capital allocators is the subject of the second.

\clearpage
\section{References}

\begin{enumerate}[leftmargin=*,labelsep=0.5em,label=\arabic*.]

\item Arslanalp, S., Eichengreen, B., Simpson-Bell, C. ``Gold as
international reserves: a barbarous relic no more?'' \emph{Journal of
International Economics} \textbf{145} (2023)
\url{https://doi.org/10.1016/j.jinteco.2023.103822}

\item Arthur, W. B. ``Competing technologies, increasing returns,
and lock-in by historical events.'' \emph{Economic Journal}
\textbf{99.394} 116--131 (1989)
\url{https://doi.org/10.2307/2234208}

\item Bernheim, B. D., Peleg, B., Whinston, M. D.
``Coalition-proof Nash equilibria I. Concepts.'' \emph{Journal of
Economic Theory} \textbf{42.1} 1--12 (1987)
\url{https://doi.org/10.1016/0022-0531(87)90099-8}

\item Biais, B., Bisiere, C., Bouvard, M., Casamatta, C. ``The
blockchain folk theorem.'' \emph{Review of Financial Studies}
\textbf{32.5} 1662--1715 (2019)
\url{https://doi.org/10.1093/rfs/hhy095}

\item BlackRock, Inc. ``iShares Bitcoin Trust ETF, Form 10-K
Annual Report for fiscal year ended 31 December 2024.''
\emph{U.S.\ Securities and Exchange Commission EDGAR}
(28 February 2025)
\url{https://www.sec.gov/cgi-bin/browse-edgar?action=getcompany&CIK=0001980994&type=10-K}

\item Brunnermeier, M. K., James, H., Landau, J.-P. ``The
digitalization of money.'' \emph{NBER Working Paper}
\textbf{26300} (2019) \url{https://doi.org/10.3386/w26300}

\item Budish, E. ``The economic limits of Bitcoin and the
blockchain.'' \emph{NBER Working Paper} \textbf{24717} (2018)
\url{https://doi.org/10.3386/w24717}

\item Chen, Y. ``A game-theoretic foundation for Bitcoin's
price.'' \emph{Working Paper} (2025).

\item Eyal, I., Sirer, E. G. ``Majority is not enough: Bitcoin
mining is vulnerable.'' \emph{Communications of the ACM}
\textbf{61.7} 95--102 (2018)
\url{https://doi.org/10.1145/3212998}

\item Farrell, J., Saloner, G. ``Standardization, compatibility,
and innovation.'' \emph{RAND Journal of Economics}
\textbf{16.1} 70--83 (1985) \url{https://doi.org/10.2307/2555589}

\item Feldstein, M. ``Inflation and the stock market.''
\emph{American Economic Review} \textbf{70.5} 839--847 (1980)
\url{https://www.jstor.org/stable/1805772}

\item Fudenberg, D., Tirole, J. \emph{Game Theory}. Cambridge:
MIT Press (1991).

\item Hirschman, A. O. \emph{Exit, Voice, and Loyalty}. Cambridge:
Harvard University Press (1970).

\item Huberman, G., Leshno, J. D., Moallemi, C. ``Monopoly without
a monopolist: an economic analysis of the Bitcoin payment
system.'' \emph{Review of Economic Studies} \textbf{88.6}
3011--3040 (2021) \url{https://doi.org/10.1093/restud/rdab014}

\item International Monetary Fund. ``World Economic Outlook
Database.'' \emph{IMF} (October 2024)
\url{https://www.imf.org/en/Publications/WEO}

\item Katz, M. L., Shapiro, C. ``Network externalities,
competition, and compatibility.'' \emph{American Economic
Review} \textbf{75.3} 424--440 (1985)
\url{https://www.jstor.org/stable/1814809}


\item Konishi, H., Ray, D. ``Coalition formation as a dynamic
process.'' \emph{Journal of Economic Theory} \textbf{110.1} 1--41
(2003) \url{https://doi.org/10.1016/S0022-0531(03)00004-8}

\item Kuran, T. \emph{Private Truths, Public Lies: The Social
Consequences of Preference Falsification}. Cambridge: Harvard
University Press (1995).

\item Menger, C. ``On the Origin of Money.'' \emph{Economic Journal}
\textbf{2.6} 239--255 (1892)
\url{https://www.jstor.org/stable/2956146}
\item Milgrom, P., Shannon, C. ``Monotone comparative statics.''
\emph{Econometrica} \textbf{62.1} 157--180 (1994)
\url{https://doi.org/10.2307/2951479}

\item Nakamoto, S. ``Bitcoin: a peer-to-peer electronic cash
system.'' \emph{White paper} (2008)
\url{https://doi.org/10.2139/ssrn.3440802}

\item Pagnotta, E. S. ``Decentralizing money: Bitcoin prices and
blockchain security.'' \emph{Review of Financial Studies}
\textbf{35.2} 866--907 (2022)
\url{https://doi.org/10.1093/rfs/hhab058}

\item Reinhart, C. M., Sbrancia, M. B. ``The liquidation of
government debt.'' \emph{Economic Policy} \textbf{30.82}
291--333 (2015) \url{https://doi.org/10.1093/epolic/eiv003}

\item Republic of El Salvador, Legislative Assembly. ``Bitcoin
Law, Legislative Decree No.\ 57.'' \emph{Diario Oficial de la
Rep\'ublica de El Salvador} (9 June 2021)
\url{https://www.diariooficial.gob.sv/diarios/do-2021/06-junio/09-06-2021.pdf}

\item Schelling, T. C. \emph{The Strategy of Conflict}. Cambridge:
Harvard University Press (1960).

\item Topkis, D. M. \emph{Supermodularity and Complementarity}.
Princeton: Princeton University Press (1998).

\item U.S.\ Securities and Exchange Commission. ``Self-Regulatory
Organizations; Order Granting Accelerated Approval of Proposed
Rule Changes to List and Trade Shares of Spot Bitcoin
Exchange-Traded Products.'' Release No.\ 34-99306
(10 January 2024)
\url{https://www.sec.gov/files/rules/sro/cboebzx/2024/34-99306.pdf}


\end{enumerate}

% ============================================================
% APPENDIX
% ============================================================

\appendix

\section{Notation}

\begin{table}[h]
\centering
\caption*{\textbf{Table A1.} Notation used throughout.}
\begin{tabularx}{\textwidth}{lX}
\toprule
Symbol & Definition \\
\midrule
$G$ & Settlement game $(N, S, u)$ \\
$N$ & Set of capital allocators \\
$S_i$ & Strategy set: $\{\text{Exit}, \text{Stay}\}$ \\
$u_i$ & Utility function \\
$a_i$ & Actor $i$'s share of total capturable capital, $\sum_i a_i = 1$ \\
$p_t$ & Aggregate capital share at Exit at time $t$: $\sum_{i \in \text{Exit}_t} a_i$ \\
$p_i^*(K, D)$ & Threshold for player $i$ at slowly varying state $(K, D)$ \\
$p_c(K, D)$ & Critical-mass threshold at $(K, D)$ \\
$K_t$ & Financial-repression recognition share, $K_t \in [0, 1]$ \\
$D_t$ & Institutional-fragmentation index, $D_t \in [0, \infty)$ \\
$\Delta_i(p; K, D)$ & Payoff differential: $u_i(\text{Exit}) - u_i(\text{Stay})$ at $(K, D)$ \\
$R_B(p)$ & Expected real return on neutral settlement asset \\
$R_F$ & Expected real return on capturable assets \\
$\sigma_B(p)$ & Volatility loading on neutral settlement asset \\
$\sigma_F$ & Volatility loading on capturable assets \\
$K_A(p)$ & Adoption penalty \\
$R_A(p)$ & Regulatory penalty (decreases as adoption normalizes compliance) \\
$K_N(p)$ & Non-adoption penalty \\
$\lambda_i, \kappa_i, \rho_i$ & Risk aversion, career penalty, regulatory penalty (all $> 0$ strictly) \\
(M1)--(M5) & Monotonicity conditions (\S3.5) \\
\bottomrule
\end{tabularx}
\end{table}

% ============================================================
% ACKNOWLEDGEMENT
% ============================================================

\paragraph*{Acknowledgement.}

The author thanks the readers who flagged the Theorem~1
ratchet-form gap in the preprint.

% ============================================================
% AUTHOR CONTRIBUTIONS
% ============================================================

\paragraph*{Author Contributions.}
S.H.\ is the sole author and conducted all aspects of the work,
including conceptualization, formal analysis, and writing.

\paragraph*{Conflict of Interest.}
The author holds bitcoin and operates bitcoingametheory.com.

\end{document}
